%% -!TEX root = AUTthesis.tex
% در این فایل، عنوان پایان‌نامه، مشخصات خود، متن تقدیمی‌، ستایش، سپاس‌گزاری و چکیده پایان‌نامه را به فارسی، وارد کنید.
% توجه داشته باشید که جدول حاوی مشخصات پروژه/پایان‌نامه/رساله و همچنین، مشخصات داخل آن، به طور خودکار، درج می‌شود.
%%%%%%%%%%%%%%%%%%%%%%%%%%%%%%%%%%%%
% دانشکده، آموزشکده و یا پژوهشکده  خود را وارد کنید
\faculty{دانشکده مهندسی کامپیوتر}
% گرایش و گروه آموزشی خود را وارد کنید
% \department{گرایش شبکه‌های کامپیوتری}
% عنوان پایان‌نامه را وارد کنید
\fatitle{
    \settextfont{B_Titr.ttf}
    پیاده‌سازی سیستم پایش شبکه‌های کامپیوتری}
% نام استاد(ان) راهنما را وارد کنید
\firstsupervisor{دکتر مسعود صبائی}
%\secondsupervisor{استاد راهنمای دوم}
% نام استاد(دان) مشاور را وارد کنید. چنانچه استاد مشاور ندارید، دستور پایین را غیرفعال کنید.
% \firstadvisor{دکتر سیاوش خرسندی}
%\secondadvisor{استاد مشاور دوم}
% نام نویسنده را وارد کنید
\name{سیدمهدی }
% نام خانوادگی نویسنده را وارد کنید
\surname{میرفندرسکی}
%%%%%%%%%%%%%%%%%%%%%%%%%%%%%%%%%%
\thesisdate{شهریور 1401}

% چکیده پایان‌نامه را وارد کنید
\fa-abstract{
چکیده در این قسمت قرار میگیرد.
}


% کلمات کلیدی پایان‌نامه را وارد کنید
\keywords{کلیدواژه اول، ...، کلیدواژه پنجم (نوشتن سه تا پنج واژه کلیدی ضروری است)}



\AUTtitle
%%%%%%%%%%%%%%%%%%%%%%%%%%%%%%%%%%
\vspace*{7cm}
\thispagestyle{empty}
\begin{center}
\includegraphics[height=5cm,width=12cm]{besm}
\end{center}